\section{Package contents, integration, and reproduction}
\label{sec:reproduction}

\subsection{Layout and source identity}

The archive is rooted at
\path{ProveIt_Unknot_Stopping_Cores_2026-10-12/}.
Its top-level README describes the recommended integration order and
the scope of each API. The important directories are:
\begin{description}[style=nextline,leftmargin=1em]
\item[\path{article/}]
This article's complete TeX sources, reference list, compiled PDF,
vector and raster figure, figure-generation script and build script.
\item[\path{integration/}]
A repository-relative patch, complete updated source files, old/new
file hashes and checked integration instructions. Existing code changes
are limited to the core producer and verifier; the saturated and
passenger APIs are added modules.
\item[\path{artifacts/column_cancellation/}]
The complete paired benchmark JSON and focused test log.
The reproducible benchmark driver and exact baseline modules are part
of the repository integration overlay.
\item[\path{artifacts/saturated/}]
The full 106-case audit, environment, source provenance, exact seven
new cases and their source certificates, the worked matrix for case 019,
focused test log and portable audit driver.
\item[\path{artifacts/validation/}]
The complete maintained-suite log, separate passenger test log and
validation summary.
\item[\path{artifacts/source/}]
The exact retained 106-case corpus, previous primitive-annulus audit,
and their provenance hashes.
\end{description}

The source integration targets repository revision
\begin{center}\texttt{704600626606a9097cfd89d4e3d2bbc63c8dd1cb}.\end{center}
A later revision
requires reviewing the affected interfaces and replaying the retained
evidence. The package does not include a duplicate repository checkout.
Its patch was checked against exact pinned base files in an isolated
tree. The complete replacements provide a convenient second way to
inspect every proposed source change.

\subsection{Apply and validate the source changes}

Let \code{PROVEIT_CHECKOUT} be the root of a checkout at the pinned
revision, and \code{RESEARCH_PACKAGE} the root of this unpacked package.
Review \path{integration/README.md} and its old/new hash manifest first.
The essential Git operations, run from the checkout root, are:
\begin{lstlisting}
git apply --check "$RESEARCH_PACKAGE/integration/integration.patch"
git apply "$RESEARCH_PACKAGE/integration/integration.patch"
\end{lstlisting}
These commands are provided for the maintainer's integration workflow.
This research run made no remote commit or repository publication.

Run the tests from
\path{Topology/UnknotRecognition/fast} in the patched checkout:
\begin{lstlisting}
python3 -m unittest tests.test_column_cancellation \
    tests.test_normal_cut_saturated \
    tests.test_raw_passenger_transport -v
python3 -m unittest discover -s tests
\end{lstlisting}
The finite producer and replay code use Python's standard library and
the existing ProveIt modules. The maintained repository's broader tests
and oracle comparisons retain their existing dependency requirements.
The retained fresh literal-cut audit used Regina 7.4.

\subsection{Reproduce the paired benchmark}

From the same \path{fast} directory:
\begin{lstlisting}
python3 benchmark_column_cancellation.py \
    --trials 5 --output reproduced-column-cancellation.json
\end{lstlisting}
The driver locates its baseline modules and corpus relative to its own
file. It compares the same source cores and verifies the successful
results before recording timings. The optional
\code{--skip-whole-cuts} flag is useful for a finite-kernel-only
experiment, but such a run cannot reproduce the complete-query row in
the article's table. The largest synthetic baseline and the complete
query make the full experiment appreciably longer than a unit test.

Timing samples are environment-dependent. Reproduction should confirm
the algebraic outcomes, source censuses and operation counters; it need
not reproduce the saved wall-clock numbers digit for digit. A proposed
optimization should retain both affected-subset and all-case timings,
including regressions.

\subsection{Reproduce the saturated source audit}

The packaged driver accepts an explicit checkout path:
\begin{lstlisting}
python3 "$RESEARCH_PACKAGE/artifacts/saturated/run_audit.py" \
    --checkout "$PROVEIT_CHECKOUT" --fresh-oracles \
    --output reproduced-saturated/audit.json
\end{lstlisting}
It reruns the source producer, disables producer entry points during
independent replay, preserves earlier successes, and invokes the
maintained Regina literal-cut oracle subject to its expansion cap.
Omitting \code{--fresh-oracles} uses the historical comparison packets
while still rerunning the new source pipeline. The output directory is
separate from the retained audit by default.

For a single retained certificate, use the source wrapper's verifier:
\begin{lstlisting}[language=Python]
import json
from fastunknot.normal_cut_saturated_verify import \
    verify_normal_cut_saturated

with open("/path/to/certificates/019.json") as source:
    case = json.load(source)
assert verify_normal_cut_saturated(
    case["triangulation"], case["coordinates"],
    case["certificate"])
\end{lstlisting}
The certificate contains the labelled source gluings. No inference from
the isomorphism signature or from the expected genus is needed.

\Needspace{13\baselineskip}
\subsection{Call the new operators}

The saturated wrapper is opt-in:
\begin{lstlisting}[language=Python]
from fastunknot.normal_cut_saturated import normal_cut_saturated
from fastunknot.normal_cut_saturated_verify import \
    verify_normal_cut_saturated

answer = normal_cut_saturated(
    triangulation, coordinates, record_certificate=True)
if answer["status"] == "COMPLETE":
    assert verify_normal_cut_saturated(
        triangulation, coordinates, answer["certificate"])
    histogram = answer["census"]["histogram"]
\end{lstlisting}
The lower-level \code{saturated_source_census} is conditional on
already authenticated geometry, base summaries and an ordered boundary
skeleton. The wrapper supplies that authentication.

For exact count transport after a complete column schedule:
\begin{lstlisting}[language=Python]
from fastunknot.core_handlebody import column_cancellation_plan
from fastunknot.raw_passenger_transport import transport_raw_passenger

counts = {0: {0: 1, 1: 2, 2: 3},
          1: {1: 1, 3: 5}, 2: {2: 1, 3: 7}}
plan = column_cancellation_plan(counts)
passenger = transport_raw_passenger(
    counts, plan["cancellations"], {0: 11, 3: 13})
assert passenger == {3: 354}
\end{lstlisting}
The original exact \code{counts} must be supplied, not a capped copy.
The operation verifies its donor schedule but does not authenticate
embedded passenger geometry.

\subsection{Build the article and figure}

From any working directory, run
\begin{lstlisting}
python3 /path/to/package/article/build.py
\end{lstlisting}
The script uses \code{latexmk} when available and otherwise runs
\code{pdflatex} three times. It needs a standard TeX installation with
the packages named in the source. All included figure files are already
built. To regenerate them from the retained benchmark JSON, run
\code{article/make_figures.py} with Matplotlib installed, then rebuild
the article. The delivered PDF was compiled, checked for unresolved
references and layout defects, rendered, and visually inspected.

\subsection{Interpretation and integration limits}

The patch upgrades raw planning and replay, and adds a broader positive
source-cut census. It does not change the recognizer's dispatcher or
its exact fallback policy. A caller must preserve \code{INCONCLUSIVE}
as an unsupported or unfinished result. A stopping-set certificate
excludes only a complete raw schedule in its named source matrix.

The next integration milestone is to preserve embedded boundary data
through these complete schedules and certify a subsequent hierarchy
transition. The present article supplies precise local theorems,
implementations and evidence on which that work can build; the global
quasi-polynomial claim remains a research objective.
